\phantomsection
\addcontentsline{toc}{chapter}{课程参考}
\markboth{课程参考}{}
\begin{center}
{\color{structurecolor}\bfseries\fontsize{14.4}{20}\selectfont 课程参考\par}
\end{center}

\noindent\textbf{课程与教师.} MIT 18.125 \textit{Measure and Integration}, 2003年秋季, 研究生课程; 教师 Jeff Viaclovsky. 公开讲义由 Ethan Brown 根据教师的课堂手写笔记排录. 以下课程信息译自官方\href{https://ocw.mit.edu/courses/18-125-measure-and-integration-fall-2003/pages/syllabus/}{教学大纲}, 核验日期为2026-10-08.

\noindent\textbf{授课时间与先修.} 每周2次课, 每次1.5小时. 官方先修为 Analysis I (18.100). 核实的页面未给出具体星期、钟点或逐次课的公历日期; 后列课历采用官方讲次.

\noindent\textbf{大纲.} 从抽象测度论与 Lebesgue 积分入门, 先建立积分及主要收敛定理, 再构造 $\R^n$ 上的 Lebesgue 测度. 大纲还列出 $L^p$ 空间、Radon--Nikodym 定理、Lebesgue 微分定理、Fubini 定理、Hausdorff 测度以及面积与余面积公式. 本册实际正文范围仍以24份讲义及书末来源对照为准: Radon--Nikodym 是编者补充, Hausdorff 测度与面积、余面积公式未纳入. 官网课程简介另提 Fourier 变换, 但这24份原件没有对应正文.

\noindent\textbf{教材.} 必读: Walter Rudin, \textit{Real and Complex Analysis}, McGraw-Hill International Editions: Mathematics Series, McGraw-Hill Education -- Europe, 1986, ISBN 9780070542341. 推荐: Frank Jones, \textit{Lebesgue Integration on Euclidean Space}, Jones \& Bartlett, 1993; Lawrence C. Evans 与 Ronald F. Gariepy, \textit{Measure Theory and Fine Properties of Function}, CRC Press, 1991. 这里只记录官方书目信息, 不归档外部教材内容.

\noindent\textbf{考核.} 官方大纲明确\textbf{没有考试}. 课程成绩以到课及提交作业为依据, 页面未给出比例. 作业页公开三组教材题号, 截止讲次分别为8、12、20; 完整索引见后页.

\clearpage
\phantomsection
\addcontentsline{toc}{chapter}{官方课历索引}
\markboth{官方课历索引}{}
\begin{center}
{\color{structurecolor}\bfseries\fontsize{14.4}{20}\selectfont 官方课历索引\par}
\end{center}
下表据\href{https://ocw.mit.edu/courses/18-125-measure-and-integration-fall-2003/pages/calendar/}{官方课历}译录讲次、主题与作业截止安排. 课历主题与实际讲义的对应并不完全一致; 本表记录课堂安排, 逐文件覆盖见书末.

\begin{longtable}{@{}p{.07\textwidth}p{.75\textwidth}p{.12\textwidth}@{}}
\toprule
讲次 & 官网主题 & 截止作业 \\
\midrule
\endfirsthead
\toprule
讲次 & 官网主题 & 截止作业 \\
\midrule
\endhead
\bottomrule
\endfoot
1 & 测度论的动机; 测度空间与 $\sigma$-代数; 可测函数的和、积、复合; Borel 集 & \\
2 & 实值可测函数及极限; 简单函数、正测度、Lebesgue 积分定义 & \\
3 & Riemann 积分; 可积与几乎处处连续; 两种积分的比较; 正测度与积分的性质 & \\
4 & 简单函数与非负函数积分可加; 单调收敛; 求和与积分交换; Fatou 引理 & \\
5 & 复函数积分、主导收敛、零测集、$\sigma$-代数的完备化 & \\
6 & $\R^n$ 上的 Lebesgue 测度; 矩形、多面集、开集、紧集; 内外测度 & \\
7 & 从有限外测度可测集到一般可测集; Lebesgue 可测域与测度空间 & \\
8 & Carath\'eodory 判据、Cantor 集、Lebesgue 可测而非 Borel 的集合 & 1 \\
9 & 平移、伸缩与旋转不变性; 不可测集 & \\
10 & 积分作为线性泛函; 正泛函 Riesz 表示; Riemann 与 Lebesgue 积分的衔接 & \\
11 & Lusin 定理; Vitali--Carath\'eodory 定理 & \\
12 & 连续函数逼近可测函数; 几乎处处收敛下的积分定理; 积分可数可加 & 2 \\
13 & Egorov 定理; 依测度与几乎处处收敛; 子列; 依测度主导收敛 & \\
14 & 凸函数、Jensen、H\"older 与 Minkowski 不等式 & \\
15 & $L^p$、赋范与 Banach 空间; Riesz--Fischer 完备性 & \\
16 & $C_c$ 在有限指数 $L^p$ 与 $C_0$ 中稠密 & \\
17 & $L^p$ 与 $\ell^p$ 包含; 局部 $L^p$; 范数凸性; 平滑函数稠密 & \\
18 & $\R^n$ 上非负函数的 Fubini 定理 & \\
19 & $\R^n$ 上 $L^1$ 函数的 Fubini 定理; 一般测度空间的乘积测度 & \\
20 & 乘积测度的 Fubini 定理、完备化与卷积 & 3 \\
21 & Young 不等式、磨光核、$C^\infty$ 在有限指数 $L^p$ 中稠密 & \\
22 & Lebesgue 积分的微积分基本定理; Vitali 覆盖; 最大函数; 弱 $L^1$ 估计 & \\
23 & Lebesgue 微分定理、Lebesgue 点与微积分基本定理第一部分 & \\
24 & 广义 Minkowski、另一 Young 证明、分布函数; 最大算子 $L^p$ 插值估计 & \\
\end{longtable}

\clearpage
\phantomsection
\addcontentsline{toc}{chapter}{官方作业安排索引}
\markboth{官方作业安排索引}{}
\begin{center}
{\color{structurecolor}\bfseries\fontsize{14.4}{20}\selectfont 官方作业安排索引\par}
\end{center}
官方\href{https://ocw.mit.edu/courses/18-125-measure-and-integration-fall-2003/pages/assignments/}{作业页}只列三组 Rudin 教材题号与截止讲次. 教材为前列1986年 \textit{Real and Complex Analysis}; 下表的章号、题号均指该教材, 不指本册章号或33道编者练习.

\begin{center}
\begin{tabular}{@{}p{.12\textwidth}p{.15\textwidth}p{.12\textwidth}p{.48\textwidth}@{}}
\toprule
作业 & 截止讲次 & 原书章号 & 官方所列题号 \\
\midrule
1 & 8 & 1 & 3--12 \\
2 & 12 & 2 & 5--9, 20 \\
3 & 20 & 3 & 4c, 4d, 10, 11, 13, 18b, 23, 25 \\
\bottomrule
\end{tabular}
\end{center}

本次核实的公开作业页没有提供完整题面、独立作业 PDF 或官方解答. 本册因此仅收入这份安排索引, \textbf{没有收入原书题面或官方解答}, 也不据题号重建题意. 正文保留的33道配套练习及解答为编者内容, 未改称官方作业. 原课程明确无考试, 因而没有需要补入的本课试卷.

这些课程信息与安排索引保留 MIT OCW、课程、教师及年份署名, 依 OCW 使用条款与 CC BY-NC-SA 4.0 使用. 教材仅作书目与题号定位, 不将 OCW 的许可扩展为外部教材题面的转载许可.
